\documentclass{article}
\usepackage{amsmath}
\usepackage{mathtools}
\usepackage{amsfonts}
\usepackage{parskip}
\newcommand*\plus{+{}}
\newcommand*\boxSizeOfMax[1]{\makebox[\widthof{Max}][c]{#1}}
\setlength\parindent{0pt}

\title{Career Development Plan Framework - DRAFT \\ [1em]
	\large Jonathan Erdmann \\ [1em]
}

\begin{document}


\section{Executive Summary}
Mission Systems conducts an annual individual development planning cycle that represents a meaningful investment of manager and 
employee time across the organization. Based on direct observation through one-on-one conversations and consistent findings across 
pulse survey data, the return on that investment is insufficient. Engineers are producing development plans that are shallow in 
quality, limited to a one-year horizon, and rarely executed against in any substantive way.

This memo proposes a structured pilot program designed to address that gap across three dimensions. The first is plan quality and executability 
— providing senior engineers with a rigorous, research-grounded framework for building a genuine career strategy rather than a compliance document. 
The second is mentorship and peer development activation — the organization contains significant relational and institutional capital at the 
senior level that is not being converted into productive developmental relationships at the scale our depth of talent should enable. The third 
is organizational development velocity — senior engineers who cannot articulate their trajectories develop more slowly, are harder to deploy 
strategically, and represent a compounding cost to succession readiness and retention over time. These three dimensions are not independent problems. 
They are expressions of the same underlying gap, and a single framework addresses all three simultaneously.

The framework proposed here is grounded in validated research across career development, goal orientation, and professional identity literature, 
and has been designed with the specific context of senior technical professionals in program-driven environments in mind. It is not a new 
administrative burden — it is a structured methodology for producing better outputs from time the organization is already committing. The author 
has applied the framework to his own career development plan, which is attached as Appendix A as a demonstration of the method and the quality of output it produces.

The pilot is structured as follows: a cohort of 10 high-potential senior engineers, selected in coordination with leadership, 
organized into five peer accountability pairs that will work collaboratively to build and pressure-test individual career strategies over 
a six-month period, anchored to the existing annual IDP cycle. Defined success criteria will be assessed at the conclusion of the pilot, and a 
findings report delivered to leadership will include a recommendation on whether and how to scale the approach across the broader organization.

The proposed next step is a 30-minute conversation to align on cohort selection, timeline, and organizational dependencies before launch, and to 
discuss the value of VP-level sponsorship in signaling the program's organizational standing to pilot participants.

\section{Problem Statement}

\subsection{Return on Administrative Duties}
The annual individual development planning cycle represents one of the more significant recurring investments of discretionary time that 
Mission Systems makes in its engineering workforce. Managers and engineers alike commit meaningful hours each cycle to a process that is, 
in principle, designed to produce clarity of direction, accelerate capability development, and strengthen the organization's talent pipeline. 
In practice, the outputs of that process fall consistently short of those objectives.

The gap is not attributable to a lack of effort or intent on the part of either managers or engineers. It is attributable to the absence of 
a framework capable of producing the quality of output the process is ostensibly designed to generate. Without that structure, the annual 
cycle defaults to what is immediately legible — near-term skill gaps, current role performance, and training requests — rather than what is 
strategically valuable: a genuine, executable career strategy grounded in a clear understanding of where the individual is going and what it 
will take to get there. The result is an organization that spends the overhead without capturing the return.

\subsection{Network and Mentorship Underutilization}

Mission Systems carries significant mentorship potential within its senior engineering ranks — technical depth, program experience, and 
professional networks that represent genuine developmental value for the next generation of senior talent. That potential is not being 
converted into productive relationships at the scale the organization's depth of expertise should enable.

The reason is structural rather than dispositional. Mentorship relationships form most naturally when both parties have a shared 
language for articulating development goals, identifying relevant experience, and defining what a productive relationship would produce. 
Without a common framework, the conditions for those conversations to begin are absent. Potential mentors have no clear mechanism for engaging,
and potential mentees have not yet developed the strategic clarity to make a specific, meaningful ask. The relationships that do form tend to 
be informal, uneven in quality, and concentrated among engineers who already possess the professional confidence to seek them out — precisely 
the engineers who need them least.

\subsection{Organizational Development Velocity}

The consequences of the gap described above are not static — they compound. An organization whose senior engineers cannot articulate their 
trajectories develops more slowly than one that can, not because the individuals are less capable, but because deliberate development is 
more efficient than incidental development. Engineers who understand where they are going make better decisions about where to invest their 
discretionary effort, which stretch assignments to seek, which relationships to build, and which capabilities to develop ahead of demand 
rather than in response to it.

At the senior engineer level — where the organization's investment per individual is highest and the lead time for developing strategic 
capability is longest — the cost of that inefficiency is most consequential. The three-to-five year horizon is where the compounding 
effect of deliberate versus incidental development becomes visible in succession readiness, capability depth, and retention. A pilot 
program that begins generating that compounding effect now produces organizational returns that a program launched two years from 
now cannot fully recover.

\section{Proposed Framework}

\subsection{Overview}

The framework proposed here is a structured methodology for building an executable career strategy. It is designed specifically for 
senior technical professionals operating in complex, program-driven environments — not adapted from a generic talent development 
instrument, but constructed from the ground up with the specific challenges of that population in mind. It converts the existing 
IDP cycle from a compliance exercise into a genuine strategic planning process by giving engineers a defined sequence of thinking, 
a common language for articulating their development goals, and a structured output artifact they own and return to over time.

The framework consists of four components that build on each other in sequence. Each component produces a specific output that becomes 
the input for the next. The terminal output of the full sequence is a Strategic Career Plan — a concise, executable document that is 
qualitatively different from a standard IDP in specificity, time horizon, and actionability. The distinction between the framework as 
a thinking tool and the plan as its output artifact is intentional: the framework is what engineers learn and apply, the plan is what 
the organization benefits from.

The framework is grounded in three bodies of established research — career development, goal orientation, and professional identity 
— selected because each addresses a dimension of the problem that the others do not. That foundation is summarized in the following section 
and developed in full in Appendix B. It is cited here not as an appeal to academic authority but because the practical design choices 
embedded in the framework derive directly from what that research has demonstrated works.

\end{document}
